\chapter{Conclusions}
\label{chap:concl}




This thesis argued that a large class of efficiency bottlenecks in BSP systems arise from
conditions unknowable before runtime, and that such conditions are systematically addressable
through adaptive mechanisms if the requisite feedback loops exist.% \minitoc% TODO% 1.% Unifying overlay topologies (CARP + intra-tier shuffling) 2.% Open table formats as a natural fit for BSP data 3.% Convergence between algebras 4.% ML-specific differences 5.% DAQs as a parallel?% 6.% DeltaFS vs ORCA: filesystems vs databases/analytics.% Open Table Formats are at the heart.% 7.% View-driven Observability% 6.0 chapter intro: thesis claim, evidence recap, OLAP × BSP discussion, chapter roadmap

\begin{itemize}
  \item In \cref{chap:carp}, we presented CARP, an adaptive range partitioner that uses
        reconstructed key distributions to drive single-pass ingestion, matching the query performance
        of a full sort while being up to $4.9\times$ faster than post-processing approaches.

  \item In \cref{chap:amr}, we presented an empirical study of placement in AMR codes, showing that
        months of diagnostic effort were driven by the inability to construct the right runtime views,
        and CPLX, a placement policy achieving up to $21.6\%$ runtime improvement once the feedback
        gap was closed by hand.

  \item In \cref{chap:orca}, we presented ORCA, a tree-based overlay network that attaches dedicated
        aggregators to a parallel application and executes SQL-based in-situ analytics over columnar
        telemetry streams, with an asynchronous control plane providing timestep-consistent semantics.
\end{itemize}

ORCA brings OLAP machinery---columnar formats, SQL planning, embeddable query
engines---into a declarative and programmable in-situ analytics framework for HPC.
Traditional in-situ techniques are fixed-function and force a choice between embedded processing
(co-located with application ranks) and offloaded processing (a separate MPI application, as in
ADIOS~\cite{Eisen2024:ADIOS}).
ORCA's execution model is a hybrid: the most valuable operators are automatically placed at source,
while heavier analytics run on a dedicated asynchronous tier.
Unlike ADIOS-style offloading, ORCA's aggregators and controller are closer to asynchronous
services than to a coupled MPI job, while still preserving timestep-consistent execution through
TS2PC.
This intersection of OLAP analytics and BSP execution turns out to be remarkably productive: the
OLAP stack provides the data model, interchange format, and query semantics, while BSP's phase
structure provides the causal boundaries that make distributed planning tractable.

The rest of this chapter reflects on this design point (\S\ref{sec:concl:retrospective}), describes
extensions within immediate reach (\S\ref{sec:concl:extensions}), and outlines longer-term
directions (\S\ref{sec:concl:future}).

\section{Steerable Observability: Retrospective and Synthesis}\label{sec:concl:retrospective}

\paragraph{Observability as view materialization}
Observability is the problem of materializing the right view of a system.
Most efficiency and correctness problems in large-scale systems are not hidden---they are obvious
once the right view exists.
A straggler is immediately visible in a per-rank breakdown of synchronization time.
A corrupted gradient is immediately visible in a distribution of weight updates.
The challenge is not analysis but construction: assembling the view from distributed fragments, at
the right granularity, fast enough to act on.
In BSP systems, views are routinely missing not because of physical limits but because of bad
abstractions: unnecessary serialization cycles, format-dictated ETL, layouts that force full scans,
tools that require re-execution to change what is collected.

\paragraph{Steerability and limits}
Steerability is what distinguishes a general observability substrate from a fixed-function monitor.
It means control over what is collected, where it is processed, how much is retained, and what is
acted on---all adjustable at runtime under explicit overhead and jitter budgets.
Steerability extends beyond observation: views inform interventions, and interventions must be
assessed through further views, closing the loop.

The concept is only viable insofar as the feedback loop latency permits.
Constructing a global view has an intrinsic cost---collection, reduction, and transport take
time---and the view is only useful if it arrives soon enough to inform a decision.
Many systems operate in regimes where this cost exceeds the decision timescale: network congestion
control algorithms, for instance, are explicitly designed around partial observability because a
complete view of path state cannot be assembled within a round-trip.
Steerable observability is productive in regimes where the feedback latency is shorter than the
timescale of the conditions it targets.
BSP workloads, with timesteps on the order of milliseconds to seconds and runs lasting weeks, sit
squarely in this regime.

\paragraph{Why BSP demands a different design}
BSP forces a specific shape on observability infrastructure that differs from cloud-native
approaches.
The distributed systems community has become microservice-brained: observability is designed around
many independent services generating loosely correlated event streams, where eventual consistency
is
acceptable and metrics can be sampled.
BSP is one workload running as a single hero run, with the fabric providing the illusion of a
single computer.
Global synchronization means every rank's behavior affects every other rank's progress.
Tail sensitivity means a single slow rank dominates an entire timestep.
Phase structure---alternating compute and I/O---imposes non-uniform overhead budgets: during
compute, any jitter from telemetry collection amplifies into synchronization cost, while during
I/O the constraints relax.
This demands reduction-first, phase-aware design: reduce at source during compute, with the option
to move more data when jitter is not the binding constraint.

\paragraph{ORCA's design point}
ORCA's design follows from these constraints.
Telemetry originates as schema-first Arrow columnar streams, avoiding the serialization overhead of
log-based formats.
Queries are SQL---a restricted but plannable operator vocabulary that enables pushdown across TBON
tiers so only the data needed for the requested view moves over RDMA.
The result is close to optimal for a given view: columnar data, reduced in-situ, no avoidable
serialization or data cycles.
The physics DAQ pipeline is the right mental model: discard and reduce at source, keep only what
the downstream question requires---but unlike a hardware trigger, ORCA's filters are steerable.
This philosophy has precedent in software: Fay pushes LINQ operators into tracepoints, and Snicket
compiles queries into WebAssembly filters inside microservice proxies.
Both validate the core move but neither operates in a BSP context with phase-aware semantics,
columnar interchange, or a timestep-consistent control plane.

\section{Immediate Extensions}\label{sec:concl:extensions}


ORCA's architecture is modular by construction: new sources produce Arrow RecordBatches, new
operators plug into DataFusion's execution framework, and new sinks attach to the overlay.
Several extensions require no architectural changes and are largely a matter of implementation
effort.

\paragraph{Instrumentation sources}
Any interface that can emit structured events into Arrow is within scope.
ORCA's \texttt{SchemaCollector} abstraction---a hot-path \texttt{AddEvent} into a typed struct,
drained to produce a RecordBatch---makes new sources mechanical to integrate, whether NVTX/CUPTI
GPU hooks, MPI runtime probes, or kernel events.
eBPF is particularly natural: dynamic probes can be attached at runtime and emit fixed-schema
events into the same pipeline.
This is immediate---eBPF as a source.
The harder problem of compiling query operators down into eBPF bytecode is a longer-term direction.

\paragraph{Approximation and sketches}
Collecting all telemetry is rarely the right default.
Most diagnostic questions are aggregate---distributions, quantiles, heavy hitters,
correlations---and exact answers over complete data are more expensive than necessary.
Mergeable sketches align naturally with ORCA's reduction topology: compute a sketch at each leaf,
merge hierarchically up the tree, emit a compact summary at the root.
This makes observability cost-proportional to the question rather than the raw event rate.
The remaining design questions are semantic: which views may be approximate, what error bounds are
acceptable, and how approximation interacts with completeness guarantees and phase alignment.

\paragraph{Streaming analytics unification}
ORCA's data plane and CARP's shuffle pipeline serve different purposes but share the same
structural pattern: planned movement of columnar data across a topology.
CARP/DeltaFS-style all-to-all shuffles fit naturally as intra-tier operators within ORCA's planner.
Extending ORCA's sink scope to include MPI-tier endpoints would allow a single plan to span
compute-time observability reduction and checkpoint-time data reorganization---different operators
activated in different phases, composed in one framework.
This collapses the current separation between observability infrastructure and data management
infrastructure.

\paragraph{Open table formats for HPC}
ORCA already writes timestep/rank-partitioned Parquet, but downstream tools see files rather than a
dataset.
Adding a metadata catalog---Iceberg-style snapshots, partition pruning, schema evolution---makes
traces and telemetry into first-class tables queryable by any engine that speaks the format.
This is a broader claim than ORCA: HPC outputs in general---simulation checkpoints, analysis
artifacts, telemetry---would benefit from open table formats that make partition structure, schema,
and versioning explicit.
ORCA's Parquet output is a natural on-ramp, but the value extends to any workflow that currently
manages collections of files by convention.

\paragraph{Control interface regularization}
ORCA's control plane currently supports timestep-consistent state transitions via TS2PC, but the
set of controllable actions is defined ad hoc per integration.
A natural step is to regularize this into typed control-plane providers that publish capabilities,
accept structured commands, and return typed acknowledgments.
This makes planned actuation composable in the same way that planned queries are
composable---the planner can reason about what control actions are available, what their
consistency
semantics are, and how they interact with active queries.

\section{Future Directions}\label{sec:concl:future}

\paragraph{Beyond relational algebra}
Relational operators are a natural fit for event-style telemetry---filtering, projection,
aggregation, joins---but they are a poor fit for the structured data that dominates scientific
computing and ML.
Arrays, meshes, and tensors have dimensionality, adjacency, and layout that determine both
semantics and efficient access patterns.
Forcing them into relational form destroys this structure.
A natural extension is a multi-algebra planner: retain relational operators for event streams, add
array and mesh operators for structured data, and compose them in one plan.
Format interop is becoming practical through Arrow bridges---Arrow Zarr for array data,
Arrow DLPack for tensors---which provide a path to querying device- and array-resident data
without premature tabularization.
Prior work on MeshSQL, SciQL, and array databases provides clear precedent for extending SQL to
cover these data types.
A motivating example: compute a histogram of ML weight values when weights live in GPU memory.
The right execution is a GPU-side reduction that emits only the compact histogram back into the
pipeline---not a transfer of the full tensor to the host for relational processing.

\paragraph{Visualization as planned dataflows}
HPC visualization workflows are typically bespoke extract-transform-render pipelines that reinvent
data movement and reduction logic for each code and toolchain.
These can be subsumed by treating visualization products---tiles, isosurfaces, decimated particles,
derived fields---as another class of sinks within the same planning framework.
The operators differ (level-of-detail construction, GPU-friendly reductions, domain-specific
kernels) but the planning and placement machinery is shared.
This unifies observability and visualization under one infrastructure, eliminating the current
duplication of data movement and reduction logic across separate systems.

\paragraph{Multi-tier compilation}
Once eBPF probes, host collectors, SmartNICs, and TBON aggregators are all potential execution
targets, a SQL query becomes a multi-target compilation and placement problem.
Each tier admits a restricted operator subset: eBPF allows bounded per-event filtering and small
sketches, DPUs support streaming transforms with limited state, higher tiers handle heavier
operators like joins and windowed aggregations.
Full SQL semantics do not port cleanly across the pipeline.
The core challenge is lowering a high-level query into tier-appropriate kernels with explicit
materialization boundaries, while preserving BSP phase semantics and meeting jitter and overhead
budgets.
This is a compiler problem as much as a systems problem.

\paragraph{Cost-aware reactive placement}
ORCA's current operator placement is largely topology-informed---per-rank data gets pushed down,
aggregation rises up the tree.
This works well in practice but is not cost-aware: it does not observe realized pipeline costs
(ingestion rates, queueing depths, CPU and NIC pressure, jitter impact) or treat placement as a
controlled variable.
A natural evolution is to add explicit cost signals and guardrails, then allow the planner to
revise placement decisions online---moving operators and sinks across tiers in response to observed
conditions rather than relying on accurate up-front prediction.
This closes a feedback loop over ORCA itself: the same infrastructure that provides steerable
observability for applications could provide steerable observability for its own execution.


% \section{Galaxy Brain Stuff}

% \subsection{Maxwell's Demon}

% It’s relevant as an analogy about “information →
% usable work via feedback,” not as a literal technical framework you should import.

% How it maps cleanly: •  Maxwell’s Demon: by measuring microstate and acting on it, the demon seems
% to extract work from thermal fluctuations.
% The resolution is that information has a cost (measurement, memory, erasure; Landauer), so you
% can’t get “free work.
% ”
% •	ORCA/CARP: by measuring distributed execution/application state and acting on it (placement, throttling, repartition, etc.), you extract performance/efficiency from variability/nonstationarity.
% The “demon tax” is observability overhead + coordination cost + perturbation, which has to be paid
% somewhere.

% So the useful “science” point is:

% Feedback can turn uncertainty/variance into
% advantage only insofar as the cost of acquiring/processing information is less than the benefit of
% better decisions.

% That’s exactly the tradeoff ORCA formalizes as “query-driven / pay proportional to the question,”
% and what CARP informally exploits with its renegotiation epochs.

% Where the analogy breaks / becomes unhelpful: • Thermodynamics is about entropy/energy with a tight
% formalism; your setting is distributed control with resource tradeoffs (bandwidth, latency, CPU,
% coordination).
% You can steal the intuition (no free lunch; info isn’t free), but don’t pretend there’s a conserved
% quantity you can compute.

% If you use it at all, it should be a 1–2 sentence motivating metaphor: “ORCA is an information
% engine; measurement + control can harvest slack, but only if measurement costs are bounded.
% ”

% \subsection{Ashby's Law of Requisite Variety}

% % remove all unicode characters
% Ashby's Law of Requisite Variety (cybernetics/control) says, roughly: a regulator can only suppress
% disturbances to the extent that the regulator has comparable "variety" (i.e., a sufficient
% repertoire of distinguishable responses / states).
% Ashby's aphorism is commonly rendered as 'only variety can destroy variety'.

% More formally (one common statement): the variety/uncertainty left in the outcomes can't be reduced
% below what remains after accounting for the regulator's variety; in information-theoretic form
% you'll see inequalities like $V_O \ge V_D - V_R$ or $H(E) \ge H(D) - H(R)$ under standard
% assumptions.

% Is it relevant to ORCA/CARP?

% Yes, as a useful lens—but you have to map “variety” correctly.
% •	In Ashby, "variety" is about the regulator's ability to select among alternatives.
% That selection capacity can come from:
% \begin{itemize}
%     \item Action variety (many possible interventions),
%     \item Information/measurement that lets you choose the right intervention conditional on the
%     \item Ashby explicitly links regulation limits to communication/channel capacity.
% \end{itemize}
% and/or Information/measurement that lets you choose the right intervention conditional on the disturbance (selection requires discrimination).

% CARP: the disturbance is (time-varying) key-distribution skew; the regulator's variety is
% essentially "many possible repartitionings + the mechanism to choose among them," based on observed
% distributions.
% It's a specialized regulator with a narrow disturbance class but a strong, domain-specific action
% set.

% ORCA: the disturbance class is broader (performance pathologies, imbalance modes, phase changes,
% cross-stack variability).
% ORCA's claim-to-power is that it expands the regulator's effective variety by: •        making the
% measurement space programmable and composable (you can form the discriminants you need), and •
%  aligning sensing/actuation to a semantic clock (logical time) so the "selection" is coherent at
% BSP boundaries (i.e., you're choosing based on a well-defined slice of the world, not an arbitrary
% asynchronous smear).

% Under Ashby's lens: ORCA is an attempt to raise $V_R$ (effective regulatory capacity) not just by
% adding knobs, but by increasing the ability to pick the right knob setting in-context—i.e.,
% improving discrimination/selection under uncertainty.
% That's conceptually aligned with Ashby even though you're not importing thermodynamic-style
% conservation laws.

% (W.Ross Ashby is the original source; his Introduction to Cybernetics is the canonical reference.)
